\chapter{最优控制与量化交易的泛函建模}

\label{chap:11}

从第 \ref{chap:01} 章到第 \ref{chap:10} 章，本书已经建立了无穷维凸优化的主要技术链：拓扑与测度负责给出合法环境，泛函分析与分离定理负责给出对偶与乘子对象，次微分、共轭与单调算子负责把最优性系统改写成可计算的原--对偶结构。到了这一章，目标不再是继续扩展抽象工具，而是把这些工具真正落到连续时间控制和量化交易的函数空间建模上。关键变化在于：决策变量不再是一个有限维向量，而是一条控制轨道、一族适应过程，或者某个随时间递推的价值函数。

本章沿三条主线展开。第 \ref{sec:11-1} 节先把连续时间最优控制写成标准的函数空间极小化问题，说明可容许控制、状态方程、成本泛函、伴随状态与一阶必要条件如何统一到第 \ref{chap:05} 章和第 \ref{chap:07} 章的语言里。第 \ref{sec:11-2} 节转向量化交易中的随机优化，把自融资策略、财富过程、风险度量与 CVaR 约束写成滤过概率空间上的凸优化模型，并用第 \ref{chap:06} 章的 Fenchel--Rockafellar 对偶与第 \ref{sec:7-4} 节的 KKT 系统解释影子价格和风险乘子。最后，第 \ref{sec:11-3} 节讨论随机控制与 Bellman 动态规划，说明价值函数、Bellman 算子和 HJB 极限怎样把第 \ref{chap:02} 章的 It\^o 语法与第 \ref{chap:10} 章的算法语义连接起来。

本章依旧坚持“应用不脱离主线”的写法。最优控制并不会被写成一部控制理论专著，量化金融也不会被扩展成一部资产定价教材。我们只保留那些会反复回到无穷维凸优化主线中的对象：直接法存在性、伴随与最优性系统、风险约束下的原--对偶结构、以及由 Bellman 递推组织起来的动态规划接口。与第 \ref{chap:08} 章有限维坍缩的写法类似，Boyd 中的 LQR、MPC、均值--方差、交易成本和执行模型在本章都只是入口；真正重要的是弄清楚这些例子在函数空间、对偶空间和过程空间中分别对应什么对象。

\section{连续时间最优控制论}

\label{sec:11-1}

连续时间最优控制最容易被有限维直觉遮蔽的地方，在于它的变量不是某个向量，而是整条控制轨道以及由状态方程生成的整条状态轨道。也正因为如此，正确的建模起点不是“写出 Hamiltonian 再求导”，而是先把控制集、状态空间、状态方程和成本泛函组织成一个真正的函数空间极小化问题，然后再讨论存在性、可微性与一阶必要条件。本节就是要把这条链条写完整。

\subsection{控制空间、状态方程与标准型}

\begin{definition}[可容许控制]\label{def:admissible-control}
设 $U$ 为 Hilbert 空间，$[0,T]$ 为给定时间区间，$K\subset U$ 为非空闭凸集。若控制函数
\[
u\in L^2(0,T;U)
\]
满足
\[
u(t)\in K \qquad \text{对几乎处处 } t\in[0,T]
\]
成立，则称 $u$ 为一个\emph{可容许控制}。所有可容许控制组成的集合记为
\[
\mathcal U_{\mathrm{ad}}.
\]
\end{definition}

\begin{definition}[控制问题的泛函标准型]\label{def:control-functional-standard-form}
设 $X$ 为状态空间，$U$ 为控制空间，$S\colon \mathcal U_{\mathrm{ad}}\to L^2(0,T;X)$ 为由状态方程生成的控制到状态映射。给定运行成本 $\ell\colon [0,T]\times X\times U\to \mathbb R\cup\{+\infty\}$ 与终端成本 $\Phi\colon X\to \mathbb R\cup\{+\infty\}$，定义
\[
J(u):=\int_0^T \ell\bigl(t,(Su)(t),u(t)\bigr)\,dt+\Phi\bigl((Su)(T)\bigr).
\]
极小化问题
\[
\min_{u\in \mathcal U_{\mathrm{ad}}} J(u)
\]
称为连续时间最优控制问题的\emph{泛函标准型}。
\end{definition}

\begin{remark}
定义~\ref{def:control-functional-standard-form} 把最优控制重新挂回了第 \ref{chap:05} 章的 proper 泛函和第 \ref{chap:07} 章的约束标准型。控制约束体现在 $\mathcal U_{\mathrm{ad}}$ 的闭凸性中，状态方程体现在映射 $S$ 中，状态约束和运行成本则都可以被吸收到 $J$ 的取值里。这样做的好处是：存在性和最优性条件可以直接调用定理~\ref{thm:direct-method-reflexive}、定义~\ref{def:gateaux-derivative} 与定义~\ref{def:kkt-cone-system} 的已有接口，而不必重新发明一套平行语法。
\end{remark}

\subsection{直接法存在性}

\begin{theorem}[最优控制问题的存在性]\label{thm:optimal-control-existence}
设 $U$ 为 Hilbert 空间，$\mathcal U_{\mathrm{ad}}\subset L^2(0,T;U)$ 为非空闭凸集。假设：
\begin{enumerate}
  \item 控制到状态映射 $S\colon \mathcal U_{\mathrm{ad}}\to L^2(0,T;X)$ 弱连续；
  \item 泛函 $J$ 在 $L^2(0,T;U)$ 上 proper、下半连续且凸；
  \item $J$ 在 $\mathcal U_{\mathrm{ad}}$ 上强制，即
  \[
  \|u_n\|_{L^2}\to\infty \quad\Longrightarrow\quad J(u_n)\to +\infty.
  \]
\end{enumerate}
则控制问题
\[
\min_{u\in \mathcal U_{\mathrm{ad}}} J(u)
\]
存在至少一个最优控制 $\bar u\in \mathcal U_{\mathrm{ad}}$。
\end{theorem}

\begin{proof}
取极小化列 $\{u_n\}\subset \mathcal U_{\mathrm{ad}}$，使
\[
J(u_n)\downarrow \inf_{u\in \mathcal U_{\mathrm{ad}}}J(u).
\]
由强制性，$\{u_n\}$ 在 $L^2(0,T;U)$ 中有界。由于 $L^2(0,T;U)$ 是 Hilbert 空间，特别是自反空间，故存在子列仍记为 $\{u_n\}$ 与某个 $\bar u\in L^2(0,T;U)$ 使
\[
u_n\rightharpoonup \bar u.
\]
又因为 $\mathcal U_{\mathrm{ad}}$ 闭凸，所以它在弱拓扑下闭，从而 $\bar u\in \mathcal U_{\mathrm{ad}}$。现在应用第 \ref{sec:5-1} 节定理~\ref{thm:direct-method-reflexive} 的直接法框架：由 $J$ 的下半连续性，
\[
J(\bar u)\le \liminf_{n\to\infty}J(u_n)=\inf_{u\in \mathcal U_{\mathrm{ad}}}J(u).
\]
故 $\bar u$ 达到极小值。
\end{proof}

\begin{remark}
定理~\ref{thm:optimal-control-existence} 的重点不是“控制问题总有最优解”，而是指出最优解来自三个具体机制：可容许控制集的弱闭性、目标泛函的弱下半连续性，以及强制性把极小化列困在一个弱紧区域中。这正是第 \ref{sec:5-1} 节直接法在控制情形中的标准落地。
\end{remark}

\subsection{伴随状态与一阶必要条件}

\begin{definition}[伴随状态]\label{def:adjoint-state-control}
设控制问题可写成
\[
J(u)=F(Su,u),
\]
其中 $F$ 在 $(x,u)$ 变量上 Fr\'echet 可微，$S$ 为有界线性算子。若对某个控制 $\bar u$，存在函数或过程 $p$ 满足
\[
S^\ast D_xF(S\bar u,\bar u)=p,
\]
则称 $p$ 为与 $\bar u$ 对应的\emph{伴随状态}。在具体微分方程模型中，它往往满足一个反向终端值问题，这正是第 \ref{sec:3-3} 节定义~\ref{def:adjoint-operator} 在控制中的具体化。
\end{definition}

\begin{theorem}[凸控制问题的 Pontryagin/KKT 型必要条件]\label{thm:pontryagin-convex-necessary}
设 $\mathcal U_{\mathrm{ad}}\subset L^2(0,T;U)$ 为非空闭凸集，$S\colon L^2(0,T;U)\to L^2(0,T;X)$ 为有界线性算子，$F\colon L^2(0,T;X)\times L^2(0,T;U)\to \mathbb R$ 为凸且 Fr\'echet 可微。若
\[
\bar u\in \operatorname*{argmin}_{u\in \mathcal U_{\mathrm{ad}}}F(Su,u),
\]
则存在伴随状态 $p$，使得对任意 $u\in \mathcal U_{\mathrm{ad}}$ 都有
\[
\left\langle S^\ast D_xF(S\bar u,\bar u)+D_uF(S\bar u,\bar u),\,u-\bar u\right\rangle \ge 0.
\]
若再把约束集写成指标函数 $\iota_{\mathcal U_{\mathrm{ad}}}$，则上式等价于
\[
0\in S^\ast D_xF(S\bar u,\bar u)+D_uF(S\bar u,\bar u)+N_{\mathcal U_{\mathrm{ad}}}(\bar u),
\]
这就是控制问题的 Pontryagin/KKT 型一阶必要条件。
\end{theorem}

\begin{proof}
定义复合泛函
\[
\mathcal J(u):=F(Su,u).
\]
由链式法则和定义~\ref{def:adjoint-state-control}，
\[
D\mathcal J(\bar u)h
=
\langle D_xF(S\bar u,\bar u),Sh\rangle
\,+\,\langle D_uF(S\bar u,\bar u),h\rangle
\]
\[
=
\left\langle S^\ast D_xF(S\bar u,\bar u)+D_uF(S\bar u,\bar u),\,h\right\rangle.
\]
由于 $\bar u$ 在闭凸集 $\mathcal U_{\mathrm{ad}}$ 上最优，应用第 \ref{sec:5-2} 节定义~\ref{def:gateaux-derivative} 所刻画的一阶变分和凸集上的最优性条件，得到对任意 $u\in \mathcal U_{\mathrm{ad}}$，
\[
D\mathcal J(\bar u)(u-\bar u)\ge 0.
\]
代入上式便得
\[
\left\langle S^\ast D_xF(S\bar u,\bar u)+D_uF(S\bar u,\bar u),\,u-\bar u\right\rangle \ge 0.
\]
再由法锥与指标函数的次微分关系，可将其改写为
\[
0\in S^\ast D_xF(S\bar u,\bar u)+D_uF(S\bar u,\bar u)+N_{\mathcal U_{\mathrm{ad}}}(\bar u).
\]
这与第 \ref{sec:7-4} 节定义~\ref{def:kkt-cone-system} 的驻点条件完全同型，只是这里的原变量是控制轨道而不是有限维向量。
\end{proof}

\subsection{典型例子}

\begin{example}[Hilbert 空间上的 LQR]
考虑线性动力系统
\[
\dot x(t)=Ax(t)+Bu(t),\qquad x(0)=x_0,
\]
以及二次成本
\[
J(u)=\frac12\int_0^T \bigl(\langle Qx(t),x(t)\rangle+\langle Ru(t),u(t)\rangle\bigr)\,dt+\frac12\langle Gx(T),x(T)\rangle.
\]
在 $Q,R,G$ 正定时，$J$ 是 Hilbert 空间上的严格凸二次泛函。把这个例子放在这里，是为了把第 \ref{sec:8-2} 节二次规划与第 \ref{sec:3-3} 节伴随算子直接桥接到控制论。
\end{example}

\begin{example}[离散时间 MPC 的有限维坍缩]
若把时间区间离散为有限网格，并把状态方程改写为
\[
x_{k+1}=A_kx_k+B_ku_k,
\]
则整个控制问题坍缩为一个有限维约束优化或二次规划问题。把这个例子放在这里，是为了说明 Boyd 中的 MPC 公式并不是另一套平行理论，而是定义~\ref{def:control-functional-standard-form} 在离散时间与欧氏空间上的直接坍缩。
\end{example}

\begin{example}[库存--产能平滑控制]
设 $x(t)$ 为库存水平，$u(t)$ 为生产速度，系统满足
\[
\dot x(t)=u(t)-d(t),
\]
并以
\[
\int_0^T \bigl(c_1x(t)^2+c_2u(t)^2\bigr)\,dt
\]
为成本。该模型既可理解为管理科学中的资源调度，也可理解为最小量级的连续时间二次控制。把这个例子放在这里，是为了给第 \ref{sec:11-2} 节交易库存惩罚模型一个确定性基线。
\end{example}

\subsection*{有限维对照}

Boyd 在有限维控制与 MPC 语境下更习惯直接写矩阵递推、Lagrangian 和 KKT 方程。本节的结论是：这些公式之所以成立，是因为离散时间与有限维把函数空间问题压缩成了标准 QP。一旦回到连续时间，真正需要先做的是定义~\ref{def:admissible-control} 与定义~\ref{def:control-functional-standard-form} 这样的对象整理，然后再用定理~\ref{thm:optimal-control-existence} 和定理~\ref{thm:pontryagin-convex-necessary} 把存在性与一阶条件接回 Boyd 的母语。有限维图像仍然重要，但它只是无穷维主线的坍缩形态。

\subsection*{习题（待补）}

\begin{exercise}[直接法存在性占位]
补一题：对一个给定的连续时间二次控制问题，验证定理~\ref{thm:optimal-control-existence} 的假设，并说明极小化列为何在弱拓扑下预紧。
\end{exercise}

\begin{exercise}[伴随系统桥接题占位]
补一题：从定理~\ref{thm:pontryagin-convex-necessary} 出发，推导一个线性二次控制模型中的伴随方程，并写出其有限维矩阵坍缩。
\end{exercise}


\section{量化金融中的随机优化}

\label{sec:11-2}

与上一节的确定性控制相比，量化交易中的真正难点并不只是噪声，而是决策本身必须适应信息流。策略是一个适应过程，财富是一个随机过程，风险约束则通常是终端损益或整个路径上的凸泛函。本节因此不把金融问题拆成“定价”和“交易”两本书，而是只保留与本书主线最贴合的部分：自融资策略、财富过程、凸风险度量、CVaR 约束，以及它们如何进入 Fenchel 对偶和 KKT 语言。

\subsection{策略、财富过程与风险泛函}

\begin{definition}[自融资策略]\label{def:self-financing-strategy}
设 $(\Omega,\mathscr F,\mathbb P)$ 为定义~\ref{def:probability-space} 的概率空间，$\{\mathscr F_t\}_{t\in[0,T]}$ 为滤过，$S=\{S_t\}_{t\in[0,T]}$ 为资产价格过程。若过程 $\pi=\{\pi_t\}_{t\in[0,T]}$ 适应于该滤过，并满足财富变化完全由持仓收益和交易成本决定，即
\[
X_t^\pi
=
x_0+\int_0^t \pi_s\,dS_s-\int_0^t c_s(\pi_s)\,ds,
\]
则称 $\pi$ 为一个\emph{自融资策略}。
\end{definition}

\begin{definition}[财富过程泛函]\label{def:wealth-process-functional}
在定义~\ref{def:self-financing-strategy} 的背景下，称
\[
\pi \longmapsto X_T^\pi
\]
为与策略 $\pi$ 对应的\emph{终端财富过程泛函}。若再给定库存或冲击惩罚 $G(\pi)$，则随机优化目标常写成
\[
\mathcal J(\pi):=\rho(-X_T^\pi)+G(\pi),
\]
其中 $\rho$ 是某个风险泛函。
\end{definition}

\begin{definition}[凸风险度量]\label{def:convex-risk-measure}
定义在 $L^1(\Omega,\mathscr F,\mathbb P)$ 上的泛函 $\rho$ 称为一个\emph{凸风险度量}，若对任意随机损失 $L,L_1,L_2$ 和 $\lambda\in[0,1]$ 满足：
\begin{enumerate}
  \item 单调性：若 $L_1\le L_2$ a.s.，则 $\rho(L_1)\le \rho(L_2)$；
  \item 平移不变性：对任意常数 $m$，
  \[
  \rho(L+m)=\rho(L)+m;
  \]
  \item 凸性：
  \[
  \rho(\lambda L_1+(1-\lambda)L_2)
  \le
  \lambda \rho(L_1)+(1-\lambda)\rho(L_2).
  \]
\end{enumerate}
\end{definition}

\begin{definition}[CVaR]\label{def:cvar-risk}
设损失随机变量 $L\in L^1$，置信水平 $\alpha\in(0,1)$。定义
\[
\operatorname{CVaR}_\alpha(L)
:=
\inf_{\eta\in\mathbb R}
\left\{
\eta+\frac{1}{1-\alpha}\mathbb E[(L-\eta)_+]
\right\}.
\]
它称为损失 $L$ 在水平 $\alpha$ 下的\emph{条件风险价值}。
\end{definition}

\subsection{CVaR 的对偶表达与风险约束}

\begin{proposition}[CVaR 的对偶表示]\label{prop:cvar-dual-representation}
对任意 $L\in L^1$ 与 $\alpha\in(0,1)$，
\[
\operatorname{CVaR}_\alpha(L)
=
\sup\left\{
\mathbb E[ZL]:
Z\ge 0,\ \mathbb E[Z]=1,\ Z\le \frac{1}{1-\alpha}\ \text{a.s.}
\right\}.
\]
\end{proposition}

\begin{proof}
由定义~\ref{def:cvar-risk}，
\[
\operatorname{CVaR}_\alpha(L)
=
\inf_{\eta\in\mathbb R}
\left\{
\eta+\frac{1}{1-\alpha}\mathbb E[(L-\eta)_+]
\right\}.
\]
对函数 $\varphi(r)=(r)_+$，其共轭满足
\[
\varphi^\ast(z)=
\begin{cases}
0, & 0\le z\le 1,\\
+\infty, & \text{otherwise}.
\end{cases}
\]
因此
\[
\frac{1}{1-\alpha}(L-\eta)_+
=
\sup_{0\le Z\le \frac{1}{1-\alpha}} Z(L-\eta)
\]
逐点成立。把该式代回并交换积分与上确界，得到
\[
\operatorname{CVaR}_\alpha(L)
=
\inf_{\eta}\sup_{0\le Z\le \frac{1}{1-\alpha}}
\left\{
\eta+\mathbb E[ZL]-\eta \mathbb E[Z]
\right\}.
\]
若 $\mathbb E[Z]\neq 1$，则右侧关于 $\eta$ 的仿射函数无下界；故只有满足 $\mathbb E[Z]=1$ 的 $Z$ 才可能给出有限值。于是
\[
\operatorname{CVaR}_\alpha(L)
=
\sup\left\{
\mathbb E[ZL]:
Z\ge 0,\ \mathbb E[Z]=1,\ Z\le \frac{1}{1-\alpha}\ \text{a.s.}
\right\}.
\]
结论成立。
\end{proof}

\begin{remark}
命题~\ref{prop:cvar-dual-representation} 说明，CVaR 并不只是一个工程上“好算”的风险指标，它本质上也是一个对偶对象。约束
\[
\operatorname{CVaR}_\alpha(L)\le \tau
\]
等价于一族有界密度权重下的最坏情形期望约束。这正是第 \ref{sec:6-4} 节定理~\ref{thm:fenchel-rockafellar-duality} 在风险建模中的典型落点。
\end{remark}

\subsection{交易执行与风险对偶}

\begin{theorem}[投资组合风险模型的对偶表达]\label{thm:portfolio-risk-duality}
设 $\mathcal A$ 为一组自融资策略构成的非空闭凸集，$\pi\mapsto X_T^\pi$ 为从 $\mathcal A$ 到 $L^1$ 的连续仿射映射，$G\colon \mathcal A\to \mathbb R\cup\{+\infty\}$ 为 proper、凸且下半连续泛函，$\rho$ 为凸风险度量。若资格条件满足，则原问题
\[
\inf_{\pi\in \mathcal A}\{\rho(-X_T^\pi)+G(\pi)\}
\]
与对偶问题
\[
\sup_{Z\in L^\infty}
\left\{
-\rho^\ast(Z)-G^\ast(K^\ast Z)
\right\},
\qquad
K\pi:=-X_T^\pi,
\]
具有相同最优值。任意最优原--对偶对 $(\bar\pi,\bar Z)$ 还满足
\[
\bar Z\in \partial \rho(-X_T^{\bar\pi}),
\qquad
0\in K^\ast \bar Z+\partial G(\bar\pi)+N_{\mathcal A}(\bar\pi).
\]
\end{theorem}

\begin{proof}
把原问题写成
\[
\inf_{\pi\in \mathcal A}\bigl\{\rho(K\pi)+G(\pi)\bigr\}.
\]
由于 $K$ 连续仿射且 $\mathcal A$ 闭凸，可以把约束通过指标函数 $\iota_{\mathcal A}$ 吸收到 $G$ 中，从而得到标准的复合凸优化模板
\[
\inf_{\pi}\bigl\{\rho(K\pi)+\widetilde G(\pi)\bigr\},
\qquad
\widetilde G:=G+\iota_{\mathcal A}.
\]
现在直接应用第 \ref{sec:6-4} 节定理~\ref{thm:fenchel-rockafellar-duality}，得到
\[
\inf_{\pi}\bigl\{\rho(K\pi)+\widetilde G(\pi)\bigr\}
=
\sup_{Z}\bigl\{-\rho^\ast(Z)-\widetilde G^\ast(K^\ast Z)\bigr\}.
\]
又因为 $\widetilde G=G+\iota_{\mathcal A}$，其最优性条件由第 \ref{sec:7-4} 节推论~\ref{cor:kkt-via-fenchel-subdifferential} 改写为
\[
\bar Z\in \partial \rho(-X_T^{\bar\pi}),
\qquad
0\in K^\ast \bar Z+\partial G(\bar\pi)+N_{\mathcal A}(\bar\pi).
\]
这正是所述原--对偶条件。对偶变量 $\bar Z$ 因而可以被解释为风险约束下的影子价格或极端情景权重。
\end{proof}

\subsection{典型例子}

\begin{example}[单期均值--方差组合作为有限维坍缩]
若交易只发生在一个时点，策略退化为有限维向量 $w$，则问题
\[
\min_w \left\{-\mu^\top w+\frac{\lambda}{2}w^\top \Sigma w\right\}
\]
正是 Boyd 书中最熟悉的均值--方差模型。把这个例子放在这里，是为了说明定义~\ref{def:self-financing-strategy} 和定义~\ref{def:wealth-process-functional} 在单期下会完全压缩成有限维投资组合优化。
\end{example}

\begin{example}[CVaR 约束下的策略优化]
若把终端损失写成 $L(\pi)$，并加入约束
\[
\operatorname{CVaR}_\alpha(L(\pi))\le \tau,
\]
则命题~\ref{prop:cvar-dual-representation} 把它改写成有界密度约束下的最坏期望控制。把这个例子放在这里，是为了把无穷维随机策略与第 \ref{chap:06} 章共轭理论真正接起来。
\end{example}

\begin{example}[带交易成本与库存惩罚的执行模型]
设 $\pi_t$ 为交易速度，$q_t$ 为库存，系统满足
\[
\dot q_t=-\pi_t,
\]
成本由冲击项 $\int_0^T c(\pi_t)\,dt$ 与库存风险项 $\int_0^T \gamma q_t^2\,dt$ 组成。该模型是量化交易执行与风险控制的标准原型。把这个例子放在这里，是为了给下一节的 Bellman 动态规划一个直接前导。
\end{example}

\begin{example}[风险中性测度作为背景例子]
在无套利定价里，对偶变量常被解释为等价鞅测度。把这个例子放在这里，只是为了提醒读者：金融中的“对偶价格”并不总是抽象泛函，它有时可以被识别为某种情景加权或测度变换；但本章主线仍然是交易执行与风险控制，而不是一般 FTAP 理论。
\end{example}

\subsection*{有限维对照}

Boyd 中的均值--方差、鲁棒优化和 CVaR 常以向量权重、样本平均和线性约束出现，因此读者很容易把它们理解成纯粹的有限维数据分析模型。本节说明，一旦把策略扩展成适应过程，真正控制问题结构的对象就变成了定义~\ref{def:self-financing-strategy} 的策略空间、定义~\ref{def:wealth-process-functional} 的财富泛函以及定理~\ref{thm:portfolio-risk-duality} 的原--对偶接口。有限维模型依旧是 Boyd 母语，但它只是过程空间优化的坍缩版本。

\subsection*{习题（待补）}

\begin{exercise}[CVaR 对偶表示占位]
补一题：从定义~\ref{def:cvar-risk} 出发，重复命题~\ref{prop:cvar-dual-representation} 的推导，并解释有界密度约束的经济含义。
\end{exercise}

\begin{exercise}[交易执行桥接题占位]
补一题：把一个带库存惩罚的交易执行模型写成定理~\ref{thm:portfolio-risk-duality} 的标准形式，并指出对偶变量对应的影子价格。
\end{exercise}


\section{随机控制与 Bellman 动态规划}

\label{sec:11-3}

上一节已经把交易执行与风险控制写成了适应过程上的凸优化问题，但这还没有回答一个更动态的问题：当决策要在每个时刻随状态和信息更新时，如何把整个问题拆成逐时递推的结构？Bellman 动态规划给出的答案是价值函数。它把“从当前状态出发的最优剩余成本”重新编码成一个函数，于是控制问题不再只是单个极小化，而是一个由 Bellman 算子驱动的递推系统。本节因此既是随机控制的结构入口，也是从 It\^o 分析走向 HJB 方程的过渡层。

\subsection{价值函数与 Bellman 算子}

\begin{definition}[随机控制的价值函数]\label{def:value-function-stochastic-control}
考虑离散时间随机控制系统
\[
X_{k+1}=F_k(X_k,u_k,\xi_{k+1}),
\qquad k=0,1,\dots,N-1,
\]
其中 $\{\xi_k\}$ 为噪声序列，控制 $u_k$ 对给定滤过适应。若阶段成本为 $\ell_k(x,u)$，终端成本为 $g(x)$，则从时刻 $k$ 和状态 $x$ 出发的\emph{价值函数}定义为
\[
V_k(x)
:=
\inf_{u_k,\dots,u_{N-1}}
\mathbb E\!\left[
\sum_{j=k}^{N-1}\ell_j(X_j,u_j)+g(X_N)
\;\middle|\;
X_k=x
\right].
\]
\end{definition}

\begin{definition}[Bellman 算子]\label{def:bellman-operator}
对任意函数 $\varphi$，定义时刻 $k$ 的 \emph{Bellman 算子}为
\[
(\mathcal T_k\varphi)(x)
:=
\inf_{u\in \mathcal U(x)}
\mathbb E\!\left[
\ell_k(x,u)+\varphi(F_k(x,u,\xi_{k+1}))
\right],
\]
其中 $\mathcal U(x)$ 是状态 $x$ 下允许的控制集合。
\end{definition}

\begin{remark}
定义~\ref{def:value-function-stochastic-control} 的关键不在于“递归求值”，而在于它把原本住在策略空间中的优化问题压缩成了状态空间上的函数问题。这样一来，第 \ref{sec:11-2} 节的自融资策略可以被视作 Bellman 递推中的控制输入，而策略空间的复杂性则被转移到价值函数和 Bellman 算子的性质分析中。
\end{remark}

\subsection{动态规划原理}

\begin{theorem}[动态规划原理]\label{thm:dynamic-programming-principle}
在定义~\ref{def:value-function-stochastic-control} 的离散时间框架下，对每个 $k=0,1,\dots,N-1$ 都有
\[
V_k=\mathcal T_k V_{k+1},
\qquad
V_N=g.
\]
换言之，
\[
V_k(x)
=
\inf_{u\in \mathcal U(x)}
\mathbb E\!\left[
\ell_k(x,u)+V_{k+1}(F_k(x,u,\xi_{k+1}))
\right].
\]
\end{theorem}

\begin{proof}
固定时刻 $k$ 与状态 $x$。对任意 admissible 策略序列 $(u_k,\dots,u_{N-1})$，先把总成本拆成第一步成本与剩余成本：
\[
\sum_{j=k}^{N-1}\ell_j(X_j,u_j)+g(X_N)
=
\ell_k(x,u_k)+\sum_{j=k+1}^{N-1}\ell_j(X_j,u_j)+g(X_N).
\]
对条件期望应用定义~\ref{def:conditional-expectation} 的塔性质，得到
\[
\mathbb E\!\left[
\sum_{j=k}^{N-1}\ell_j(X_j,u_j)+g(X_N)\mid X_k=x
\right]
\]
\[
=
\mathbb E\!\left[
\ell_k(x,u_k)+
\mathbb E\!\left[
\sum_{j=k+1}^{N-1}\ell_j(X_j,u_j)+g(X_N)\mid X_{k+1}
\right]
\middle|\; X_k=x
\right].
\]
内层条件期望按定义恰好不小于 $V_{k+1}(X_{k+1})$，于是对任意第一步控制 $u$ 都有
\[
V_k(x)\ge
\inf_{u\in \mathcal U(x)}
\mathbb E\!\left[
\ell_k(x,u)+V_{k+1}(F_k(x,u,\xi_{k+1}))
\right].
\]
反过来，对任意 $\varepsilon>0$，可为下一阶段选取 $\varepsilon$-最优策略，使剩余成本不超过 $V_{k+1}$ 加上 $\varepsilon$。再把第一步控制与该近似最优尾部策略拼接，就得到反向不等式。令 $\varepsilon\downarrow 0$ 即得
\[
V_k=\mathcal T_kV_{k+1}.
\]
\end{proof}

\subsection{Bellman 迭代下的凸性保持}

\begin{proposition}[Bellman 算子的凸性保持]\label{prop:bellman-convexity-preservation}
假设状态动力学对 $(x,u)$ 仿射，控制集 $\mathcal U(x)$ 在相应意义下凸，阶段成本 $\ell_k(x,u)$ 关于 $(x,u)$ 凸且下半连续，且 $\varphi$ 为凸下半连续函数。则 Bellman 算子 $\mathcal T_k$ 仍将 $\varphi$ 送到一个凸下半连续函数：
\[
\varphi\ \text{凸下半连续}
\quad\Longrightarrow\quad
\mathcal T_k\varphi\ \text{凸下半连续}.
\]
\end{proposition}

\begin{proof}
固定任意两个状态 $x_1,x_2$ 与 $\lambda\in[0,1]$，记
\[
x_\lambda:=\lambda x_1+(1-\lambda)x_2.
\]
任取 $u_i\in \mathcal U(x_i)$，并令
\[
u_\lambda:=\lambda u_1+(1-\lambda)u_2\in \mathcal U(x_\lambda).
\]
由动力学的仿射性，
\[
F_k(x_\lambda,u_\lambda,\xi)
=
\lambda F_k(x_1,u_1,\xi)+(1-\lambda)F_k(x_2,u_2,\xi).
\]
再由 $\ell_k$ 与 $\varphi$ 的凸性，
\[
\ell_k(x_\lambda,u_\lambda)+\varphi(F_k(x_\lambda,u_\lambda,\xi))
\]
\[
\le
\lambda\bigl[\ell_k(x_1,u_1)+\varphi(F_k(x_1,u_1,\xi))\bigr]
+
(1-\lambda)\bigl[\ell_k(x_2,u_2)+\varphi(F_k(x_2,u_2,\xi))\bigr].
\]
对噪声取期望，并在两边分别对控制取下确界，即得
\[
(\mathcal T_k\varphi)(x_\lambda)
\le
\lambda (\mathcal T_k\varphi)(x_1)+(1-\lambda)(\mathcal T_k\varphi)(x_2).
\]
因此 $\mathcal T_k\varphi$ 凸。下半连续性由期望的线性、积分型下半连续性以及对控制取下确界的稳定性得到。
\end{proof}

\begin{remark}[HJB 极限接口]\label{remark:hjb-limit-control}
若把离散时间步长缩小到 $\Delta t$，并将状态递推写成
\[
X_{k+1}=X_k+b(X_k,u_k)\Delta t+\sigma(X_k,u_k)\Delta W_k,
\]
则动态规划原理配合定理~\ref{thm:ito-formula} 的 It\^o 展开，在形式上导出连续时间 HJB 方程
\[
-\partial_t V(t,x)
=
\inf_u
\left\{
\ell(x,u)+\langle b(x,u),\nabla_x V(t,x)\rangle
+
\frac12\operatorname{Tr}\bigl(\sigma\sigma^\top(x,u)\nabla_x^2V(t,x)\bigr)
\right\}.
\]
这一极限接口是随机控制和 PDE 之间最重要的桥，但本书不在此展开完整的 viscosity solution 理论；本节只强调它与 Bellman 递推的结构联系。
\end{remark}

\subsection{典型例子}

\begin{example}[最优执行的离散 Bellman 递推]
设状态为当前库存 $q_k$ 和价格因子 $Y_k$，控制为交易速度 $u_k$，并满足
\[
q_{k+1}=q_k-u_k.
\]
若阶段成本由冲击项、库存风险项和成交偏差组成，则价值函数满足一个显式的 Bellman 递推。把这个例子放在这里，是为了说明第 \ref{sec:11-2} 节的交易执行模型一旦按时间展开，就自然进入定义~\ref{def:bellman-operator} 的框架。
\end{example}

\begin{example}[连续时间二次执行模型的 HJB 影子]
在 Almgren--Chriss 型连续时间执行模型中，价值函数通常满足一个二次结构的 HJB 方程。把这个例子放在这里，是为了说明 remark~\ref{remark:hjb-limit-control} 不是形式游戏，而是交易执行与随机控制之间的真实接口。
\end{example}

\begin{example}[受风险约束的随机控制]
若在 Bellman 递推中额外引入
\[
\operatorname{CVaR}_\alpha
\]
或其他凸风险度量，就会得到一个扩展状态变量或拉格朗日乘子的动态规划系统。把这个例子放在这里，是为了把第 \ref{sec:11-2} 节的风险约束与本节的价值函数方法真正闭环。
\end{example}

\subsection*{有限维对照}

Boyd 更常从有限维动态规划、LQR 递推或最短路式 Bellman 方程讲起，因此状态和控制都默认是向量。本节说明，这些写法依旧是正确的，但它们只是定义~\ref{def:value-function-stochastic-control} 和定理~\ref{thm:dynamic-programming-principle} 在有限维 Markov 控制中的坍缩版本。真正重要的是：Bellman 方法把策略空间极小化重新编码成状态空间上的函数递推，而这一点在无穷维控制、过程控制和量化交易中同样成立。与第 \ref{sec:10-1} 节前向--后向分裂一样，它的核心不是公式多漂亮，而是结构能否稳定传递。

\subsection*{习题（待补）}

\begin{exercise}[Bellman 递推占位]
补一题：对一个有限时域最优清算模型写出价值函数和 Bellman 算子，并从定理~\ref{thm:dynamic-programming-principle} 推导递推公式。
\end{exercise}

\begin{exercise}[HJB 接口桥接题占位]
补一题：从离散时间执行模型出发，解释 remark~\ref{remark:hjb-limit-control} 中 HJB 方程的 formal limit 是如何出现的，并指出其中用到了哪一步 It\^o 语法。
\end{exercise}

